% ============================================================================
% BOLUM 4: MATERYAL VE YONTEM
% ============================================================================
% Kurumsal Temizlik Takip ve Yonetim Sistemi (KTTYS)
% Muhammet Anil YAGIZ - 213255046
% Kirikkale Universitesi - Bilgisayar Muhendisligi Bolumu
% ============================================================================

\chapter{MATERYAL VE YONTEM}
\label{chap:materyal-yontem}

\section{Giris}
\label{sec:materyal-giris}

Bu bolumde, Kurumsal Temizlik Takip ve Yonetim Sisteminin (KTTYS) 
gelistirilmesinde kullanilan programlama dilleri, framework'ler, 
veritabani sistemleri, gelistirme araclari ve metodolojiler detayli 
olarak aciklanmaktadir.

\section{Programlama Dili}
\label{sec:programlama-dili}

KTTYS projesinde, hem backend hem de frontend gelistirmesi icin 
\textbf{JavaScript} programlama dili tercih edilmistir.

\subsection{JavaScript Dili Ozellikleri}
\label{subsec:javascript-ozellikleri}

JavaScript, 1995 yilinda Brendan Eich tarafindan gelistirilen, 
yuksek seviyeli, dinamik tipli ve prototip tabanli bir programlama 
dilidir.

JavaScript'in temel ozellikleri sunlardir:

\begin{enumerate}
    \item \textbf{Dinamik Tipleme:} Degisken tipleri calisma 
    zamaninda belirlenir.
    \item \textbf{Birinci Sinif Fonksiyonlar:} Fonksiyonlar 
    degiskenlere atanabilir, parametre olarak gecilebilir.
    \item \textbf{Prototip Tabanli Kalitim:} Sinif tabanli 
    kalitim yerine, nesneler dogrudan miras alabilir.
    \item \textbf{Olay Gudumlu Programlama:} Kullanici 
    etkilesimlerine ve asenkron olaylara yanit verebilme.
    \item \textbf{Asenkron Programlama:} Promise ve async/await 
    yapilariyla asenkron islemlerin yonetimi.
\end{enumerate}

\subsection{JavaScript Secim Gerekcesi}
\label{subsec:javascript-gerekce}

JavaScript'in proje icin tercih edilme nedenleri:

\begin{itemize}
    \item \textbf{Tam Yigin Gelistirme:} Node.js sayesinde hem 
    sunucu hem de istemci tarafinda ayni dilin kullanilmasi.
    \item \textbf{Genis Ekosistem:} npm uzerinde 1.5 milyondan 
    fazla paket bulunmaktadir.
    \item \textbf{Tarayici Destegi:} Tum modern web tarayicilari 
    tarafindan yerel olarak desteklenmektedir.
    \item \textbf{Topluluk Destegi:} Dunyanin en populer 
    programlama dillerinden biri.
\end{itemize}

\section{Programlama Platformu}
\label{sec:programlama-platformu}

KTTYS'nin backend gelistirmesinde \textbf{Node.js} calisma ortami 
ve \textbf{Express.js} framework'u kullanilmistir.

\subsection{Node.js Calisma Ortami}
\label{subsec:nodejs}

Node.js, JavaScript'in sunucu tarafinda calismasini saglayan, 
acik kaynakli, platformlar arasi bir calisma ortamidir. Ryan Dahl 
tarafindan 2009 yilinda gelistirilmistir.

Node.js'in temel ozellikleri:

\begin{enumerate}
    \item \textbf{Olay Dongusu (Event Loop):} Tek is parcacikli, 
    engellemesiz (non-blocking) I/O modeli.
    \item \textbf{npm:} Dunyanin en buyuk yazilim kayit defteri.
    \item \textbf{Aktif Gelistirme:} Surekli guncellenen kararli 
    ve LTS surumleri.
    \item \textbf{Platformlar Arasi:} Windows, macOS ve Linux 
    uzerinde calisabilme.
\end{enumerate}

\subsection{Express.js Framework}
\label{subsec:expressjs}

Express.js, Node.js icin minimal ve esnek bir web uygulama 
framework'udur.

Express.js'in proje kapsaminda kullanilan ozellikleri:

\begin{enumerate}
    \item \textbf{Yonlendirme (Routing):} HTTP isteklerinin 
    uygun isleyicilere yonlendirilmesi.
    \item \textbf{Ara Katman (Middleware):} Istek-yanit dongusune 
    eklenen fonksiyonlar zinciri.
    \item \textbf{Statik Dosya Sunumu:} HTML, CSS, JavaScript ve 
    gorsel dosyalarin sunulmasi.
    \item \textbf{Hata Yonetimi:} Merkezi hata yakalama ve 
    isleme mekanizmasi.
\end{enumerate}

Express.js ile temel sunucu yapilandirmasi:

\begin{lstlisting}[language=JavaScript, caption={Express.js Sunucu Yapilandirmasi}]
const express = require('express');
const session = require('express-session');
const cors = require('cors');
const path = require('path');

const app = express();
const PORT = process.env.PORT || 3000;

// Middleware yapilandirmasi
app.use(cors({ origin: true, credentials: true }));
app.use(express.json());
app.use(express.urlencoded({ extended: true }));

// Oturum yonetimi
app.use(session({
    secret: 'kirikkale-temizlik-secret-2024',
    resave: false,
    saveUninitialized: false,
    cookie: {
        secure: process.env.NODE_ENV === 'production',
        httpOnly: true,
        maxAge: 24 * 60 * 60 * 1000 // 24 saat
    }
}));

// Sunucu baslatma
app.listen(PORT, () => {
    console.log(`Server calisiyor: http://localhost:${PORT}`);
});
\end{lstlisting}

\section{Veritabani}
\label{sec:veritabani}

KTTYS projesinde veritabani yonetim sistemi olarak \textbf{SQLite} 
tercih edilmistir.

\subsection{SQLite Ozellikleri}
\label{subsec:sqlite-ozellikleri}

SQLite, sunucu gerektirmeyen, yapilandirma gerektirmeyen, 
kendi kendine yeten bir SQL veritabani motorudur.

SQLite'in temel ozellikleri:

\begin{enumerate}
    \item \textbf{Sunucusuz Mimari:} Ayri bir veritabani sunucusu 
    gerektirmez.
    \item \textbf{Sifir Yapilandirma:} Kurulum veya yonetim 
    gerektirmez.
    \item \textbf{Tasinabilirlik:} Veritabani tek bir dosyada 
    saklanir.
    \item \textbf{ACID Uyumlu:} Atomicity, Consistency, Isolation, 
    Durability ozelliklerini destekler.
    \item \textbf{Hafif:} Yaklasik 600 KB kod boyutu.
\end{enumerate}

\subsection{sql.js Kutuphanesi}
\label{subsec:sqljs}

Node.js ortaminda SQLite erisimi icin \textbf{sql.js} kutuphanesi 
kullanilmistir. sql.js, SQLite'in JavaScript'e derlenmis 
versiyonudur.

\section{Kutuphaneler}
\label{sec:kutuphaneler}

Projede kullanilan baslica npm paketleri ve islevleri Tablo 
\ref{tab:npm-paketleri}'de listelenmistir.

\begin{longtable}{|p{3cm}|p{2cm}|p{7cm}|}
\hline
\textbf{Paket Adi} & \textbf{Versiyon} & \textbf{Islev} \\
\hline
\endfirsthead
\hline
\textbf{Paket Adi} & \textbf{Versiyon} & \textbf{Islev} \\
\hline
\endhead
\hline
\endfoot

express & 4.18.2 & Web uygulama framework'u \\
\hline
express-session & 1.17.3 & Sunucu tarafli oturum yonetimi \\
\hline
sql.js & 1.x & JavaScript SQLite veritabani motoru \\
\hline
bcryptjs & 2.4.3 & Sifre hashleme ve dogrulama \\
\hline
cors & 2.8.5 & Cross-Origin Resource Sharing yonetimi \\
\hline

\caption{Projede Kullanilan npm Paketleri}
\label{tab:npm-paketleri}
\end{longtable}

\section{Gelistirme Araclari}
\label{sec:gelistirme-araclari}

\subsection{Visual Studio Code}
\label{subsec:vscode}

Visual Studio Code (VS Code), Microsoft tarafindan gelistirilen, 
ucretsiz ve acik kaynakli bir kod editorudur.

Projede kullanilan VS Code eklentileri:
\begin{itemize}
    \item ESLint: JavaScript kod kalitesi ve stil kontrolu
    \item Prettier: Otomatik kod formatlama
    \item REST Client: API test istekleri
    \item SQLite Viewer: Veritabani icerigi goruntuleme
\end{itemize}

\subsection{Git Versiyon Kontrolu}
\label{subsec:git}

Git, dagitik versiyon kontrol sistemi olarak kullanilmistir.

\section{Yazilim Gelistirme Metodolojisi}
\label{sec:metodoloji}

Projede \textbf{Cevik (Agile)} yazilim gelistirme metodolojisi 
benimsenmistir.

\subsection{Gelistirme Sureci}
\label{subsec:gelistirme-sureci}

Proje gelistirme sureci asagidaki asamalardan olusmaktuntur:

\begin{enumerate}
    \item \textbf{Gereksinim Analizi:} Problem tanimi, paydas 
    analizi ve fonksiyonel gereksinimlerin belirlenmesi.
    \item \textbf{Tasarim:} Sistem mimarisi, veritabani semasi 
    ve API tasariminin yapilmasi.
    \item \textbf{Gelistirme:} Iteratif kod gelistirme, modul 
    bazli implementasyon.
    \item \textbf{Test:} Birim testleri, entegrasyon testleri.
    \item \textbf{Dagitim:} Sistem kurulumu ve kullanici egitimi.
\end{enumerate}

\subsection{MVC Mimari Deseni}
\label{subsec:mvc-deseni}

KTTYS, Model-View-Controller (MVC) mimari desenini takip 
etmektedir. Bu desen, uygulamayi uc temel bilesene ayirarak 
kodun organizasyonunu ve bakimini kolaylastirmaktadir.

\begin{figure}[H]
\centering
\begin{tikzpicture}[
    node distance=2cm,
    box/.style={rectangle, draw, fill=blue!20, text width=3cm, 
                text centered, minimum height=1.2cm, rounded corners}
]
    % MVC Bilesenleri
    \node[box, fill=green!30] (view) {View\\(public/)};
    \node[box, fill=orange!30, right=of view] (controller) {Controller\\(routes/)};
    \node[box, fill=purple!30, right=of controller] (model) {Model\\(database/)};
    
    % Kullanici
    \node[above=1cm of view] (user) {Kullanici};
    
    % Oklar
    \draw[->] (user) -- (view);
    \draw[<->] (view) -- (controller);
    \draw[<->] (controller) -- (model);
\end{tikzpicture}
\caption{KTTYS MVC Mimarisi}
\label{fig:mvc-mimari}
\end{figure}

\section{Proje Dosya Yapisi}
\label{sec:dosya-yapisi}

KTTYS projesinin dosya ve dizin yapisi asagidaki gibi 
organize edilmistir:

\begin{lstlisting}[basicstyle=\ttfamily\small, frame=single, 
                   caption={Proje Dizin Yapisi}]
temizliktakip/
|-- data/
|   +-- kirikkale.db      # SQLite veritabani
|-- database/
|   |-- db.js             # Veritabani baglantisi
|   |-- init.js           # Sema olusturma
|   |-- schema.sql        # SQL sema tanimlari
|   +-- seed.js           # Ornek veri ekleme
|-- middleware/
|   +-- auth.js           # Kimlik dogrulama
|-- public/
|   |-- css/style.css     # CSS stilleri
|   |-- js/app.js         # Uygulama mantigi
|   +-- index.html        # Ana HTML sayfasi
|-- routes/
|   |-- auth.js           # Kimlik dogrulama
|   |-- faculties.js      # Fakulte yonetimi
|   |-- users.js          # Kullanici yonetimi
|   |-- tasks.js          # Gorev yonetimi
|   +-- schedules.js      # Program yonetimi
|-- server.js             # Ana sunucu dosyasi
+-- package.json          # Proje bagimliliklari
\end{lstlisting}
